\documentclass[11pt,a4paper]{article}
\usepackage[margin=24mm]{geometry}
\usepackage{amsmath,amssymb,booktabs,array}
\usepackage[hidelinks]{hyperref}
\setlength{\parindent}{0pt}
\setlength{\parskip}{6pt}
\emergencystretch=2em
\newcommand{\TST}{\operatorname{TST}}
\begin{document}
\begin{center}
{\Large BP-R Stage 3: accounting and reference contract}\\[4pt]
6 October 2026 --- reference and upstream-admission A selected; source laws pending
\end{center}

\textbf{Purpose.} Explain the shared conservative accounting and the difference
between traffic quantity and disturbance-event metadata. These identities do not
adopt a traffic law, prove control leverage or validate gradients. The proposed
task has one local controlled ramp and the same open-loop rate permissions in
every independently evolving model.

\section*{1. Count each part of the population once}
Let $\mathcal C$ be the disjoint road/ramp cells or lane segments, with $n_c(t)$
their continuous vehicle-equivalent inventories. Let $q_e(t)$ denote waiting
upstream of a full modeled entrance $e$. Physical ramp approach traffic belongs
to $\mathcal C$; do not also count it as external waiting. If a point queue is
used instead of a physical approach, it replaces that approach inventory.

For offered external arrival rate $\lambda_e$, actual entrance rate $b_e$, sink
exit rate $z_c$, and conservative internal transfer rate $f_{cd}$:
\begin{align}
\dot q_e &= \lambda_e-b_e, \label{queue}\\
\dot n_c &= b_c+\sum_{d} f_{dc}-\sum_d f_{cd}-z_c. \label{cells}
\end{align}
Here $b_c$ sums external admissions into cell $c$; only the actual sink cells
have $z_c>0$. Longitudinal, lateral and merge transfers are all internal $f$'s.
Their debits and credits cancel in the population sum. Storage and receiving
restrict admissible flows; they do not delete unserved offered demand.

\begin{center}
\begin{tabular}{ll}
\toprule
Quantity & Units\\\midrule
$n_c,q_e,M,N^{\rm arr},N^{\rm exit}$ & vehicle equivalents\\
$\lambda_e,b_e,f_{cd},z_c,r$ & vehicle equivalents/second\\
Density, speed, acceleration & veh/m, m/s, m/s$^2$\\
$\TST$ & vehicle-equivalent seconds\\\bottomrule
\end{tabular}
\end{center}

Define total present population and cumulative arrivals/exits:
\begin{align}
M(t)&=\sum_c n_c(t)+\sum_e q_e(t),\\
\dot M(t)&=\Lambda(t)-Z(t),\qquad
\Lambda=\sum_e\lambda_e,\quad Z=\sum_c z_c,\\
M(t)&=M_0+N^{\rm arr}(t)-N^{\rm exit}(t). \label{mass}
\end{align}
The sums involve actual continuous traffic mass, not the number of disturbance
descriptors. Counting a descriptor as extra traffic after its mass is already
in the fluid would violate \eqref{mass}.

\newpage
\section*{2. Complete system time is the area of present traffic}
For a finite offered population, with arrivals ending at $T_a$, define
\begin{equation}
M_{\rm pop}=M_0+\int_0^{T_a}\Lambda(t)\,dt,\qquad
\TST=\int_0^{\infty} M(t)\,dt,\qquad
J=\TST/M_{\rm pop}. \label{tst}
\end{equation}
$J$ is mean remaining system time from the declared start, including waiting
before admission. Initially present traffic contributes its remaining time;
unobserved pre-start time is not reconstructed. Normalization is optional and
does not change action rankings when $M_{\rm pop}$ is common.

Why \eqref{tst}? Each infinitesimal unit of mass contributes one unit of inventory
at every instant between its offered arrival and final exit. Integrating these
contributions over time counts exactly its system time. Summing first over
traffic or first over time gives the same nonnegative area. No additional merge
work or recovery penalty is part of this objective: a physical disturbance
affects \emph{when traffic exits}, hence the inventory area.

Integrating \eqref{mass} on a finite horizon and changing integration order gives
\begin{equation}
\TST_T=M_0T+\int_0^T(T-s)\,[\Lambda(s)-Z(s)]\,ds. \label{weighted}
\end{equation}
The weights $T-s$ show why total exits at the endpoint are insufficient. For two
actions $r_a,r_b$ with the same initial mass and offered arrivals:
\begin{equation}
\TST_T[r_b]-\TST_T[r_a]
=-\int_0^T\bigl(N^{\rm exit}_{b}(t)-N^{\rm exit}_{a}(t)\bigr)\,dt. \label{delta}
\end{equation}
If the full inventory integrals are finite, take the limit $T\to\infty$.
Identical full exit histories imply identical TST even if waiting shifts between
the ramp, freeway and entrance. Equal eventual exit totals do not imply equal
exit histories or equal TST.

\textbf{Accounting example, not a freeway witness.} Ten units present initially,
all exiting at 60 s or at 120 s, give the same eventual exit count but respectively
600 and 1200 vehicle-equivalent seconds. This is an area identity, not evidence
of a feasible ramp mechanism or a claimed control benefit.

\textbf{Clearing tail.} After the decision horizon, every model uses the same
declared draining policy and remaining forcing. If exact emptying occurs, the
integral ends there. Some conservative discretizations empty only asymptotically:
then bound the omitted tail instead of requiring a fabricated exact-zero state.
For example, a justified bound $M(t)\le B\exp[-\gamma(t-T)]$ for $t\ge T$,
with $\gamma>0$, implies omitted area at most $B/\gamma$. This is a conditional
error bound, \emph{not} a proposed exponential traffic or recovery law.
Finite demand alone does not prove clearance if some traffic can remain blocked.

\newpage
\section*{3. Rate permission and disturbance creation are different maps}
A permission $r(t)$ now limits admission to the finite KW ramp approach under
approved S3-G2-M.A. It is not the realized ramp entry: propagation and receiving
intervene. In a common time-bin accounting convention let $q_k$ be prior
waiting mass, $a_k$ the offered mass available in the bin, $R^a_k$ approach
receiving in veh/s and $\Delta t$ its duration. The work-conserving entrance
interface admits continuous mass $B_k$:
\begin{equation}
B_k=\min\{q_k+a_k,\ r_k\Delta t,\ R^a_k\Delta t\},\qquad
q_{k+1}=q_k+a_k-B_k. \label{admission}
\end{equation}
The approach receives exactly $B_k$; no mass is rounded. This transcribes the
selected cap and availability contract, not a red/green launch law or the
source's full numerical solver. Receiving uses the declared current state;
the arrival-bin convention must be common to all models.
At a point merge, write the joint admissible allocation abstractly as
\begin{equation}
(f_m,f_r)=\mathcal A(D_m,D_r,R;X),\qquad
0\le f_m\le D_m,\quad 0\le f_r\le D_r,\quad f_m+f_r\le R. \label{node}
\end{equation}
$X$ is the declared current physical state; $D_m,D_r$ are demands and $R$ is
shared receiving rate. The selected upstream cap applies to entrance release,
not automatically to delayed merge entry. Stored ramp traffic can enter above
the contemporaneous entrance permission. No allocation formula,
priority fraction or capacity reduction is selected by \eqref{node}.

In the examined hybrid lineage, continuous lane-transfer intensities generate
discrete accelerating obstruction events. For one movement, an illustrative
deterministic cumulative convention has
\begin{equation}
\eta(t)=\eta_0+\int_0^t\phi(s)\,ds,\qquad
K(t)=\lfloor\eta(t)\rfloor-\lfloor\eta_0\rfloor. \label{events}
\end{equation}
$\phi$ is actual lane-transfer intensity integrated over the movement's spatial
support, in veh/s; $\eta$ is cumulative vehicle-equivalent transfer progress and
$K$ counts unit-threshold disturbance births. At a point merge the relevant
transfer is actual merge entry. The phase $\eta_0$ must be public if this convention
is used. Equation \eqref{events} illustrates the working paper's deterministic
option; it is \emph{not an adopted ramp actuator or a replacement conservation law}.

A disturbance descriptor can retain target lane, birth time, position, speed
and bounded-acceleration history. Its coupling constrains the surrounding flow.
Creating one descriptor is not an instruction to replace \eqref{cells} with
integer transfers. The source's exact mass ownership and coupling to moving
constraints still require qualification. A fractional remainder cannot
simply disappear, and fluid mass cannot be counted twice.

\textbf{Sensitivity implication only.} Event times can move when rate changes
move cumulative crossings, even while event count stays fixed. At crossings
and changes of constraint topology, an ordinary smooth derivative is not assured.
Differentiating sampled floor values alone may miss event-time effects. This
flags Stage 4 work; no derivative formula, gradient estimator, smoothing or
usefulness result is adopted here.

\textbf{Sources and transfer status.} The continuous conservation and event lead
comes from Laval--Daganzo's
\href{https://escholarship.org/content/qt8w70q261/qt8w70q261.pdf}{2004 working paper,
Sections 2--3}; the local merge lead is
\href{https://bpb-us-e1.wpmucdn.com/sites.gatech.edu/dist/f/3571/files/2023/03/merges.pdf}{Laval--Cassidy--Daganzo's merge manuscript, Sections 2--4}.
\href{https://horowitz.me.berkeley.edu/Publications_files/All_papers_numbered/80J_Gomes_TRC06.pdf}{Gomes--Horowitz 2006, Section 4}
motivates the clearance qualification. Equations \eqref{queue}--\eqref{delta}
are reviewer bookkeeping deductions under the inherited contract; the finite
task and reference transfer have not been empirically validated.
\end{document}
